\section{\parmesan}
\label{parmesan}
%\addcontentsline{toc}{section}{Parmesan}

%In preceding chapters, we have highlighted two key observations. Firstly, the adoption of a~straightforward greedy paraphrasing approach yields a substantial improvement in the correlation between BLEU and human judgment. Secondly, the Czech Meteor tables, due to their excessive noise,  have the potential to compromise the correlation of BLEU with human judgment when utilized for paraphrasing.

%This raises a fundamental questions - can the noisy nature of the Meteor tables detrimentally impact the performance of the metric? And furthermore - our methods is independent of the evaluation metric used, so why not to experiment with other metrics than BLEU?

%We kept the best performing method from the previous section \emph{one-word only} with \emph{full Depfix} and evaluate it using three common evaluation metrics -- BLEU, Meteor and TER. Based on the results (see \Tref{choosing}), we decided to employ Meteor as our metric because it shows consistently highest correlations.


In the preceding section, we demonstrated that targeted paraphrasing of reference sentences substantially improves the correlation between automatic evaluation metrics and human judgments when using simple string-based metrics such as \bleu. However, it remains unclear whether this improvement extends to evaluation metrics that already incorporate paraphrase matching internally.

\begin{mdframed}[
    linecolor=gray,
    linewidth=2pt,
    roundcorner=5,
    backgroundcolor=gray!5
]
\noindent This naturally raises two follow-up research questions: 

\begin{itemize}
    \item Does the improvement in correlation achieved by targeted paraphrasing of reference sentences for \bleu also generalize to other commonly used evaluation metrics, in particular \mn, which already incorporates paraphrase matching?
    \item Does the inherent noise in \cmpt actually degrade the performance of the \mn metric itself?
\end{itemize}

\end{mdframed}


\begin{table}[b]
\centering
% \small
\begin{tabular}{cccc}
\toprule
Metric  & Reference &  \textbf{WMT12}  &  \textbf{WMT13} \\
\midrule
\multirow{2}{*}{\bleu}   & Original & 0.751 & 0.835 \\
                        %\cline{2-4}
                        & Paraphrased  & 0.834 & \textbf{0.891} \\
\midrule
\multirow{2}{*}{\mn}    & Original & 0.833 & 0.817 \\
                        %\cline{2-4}
                        & Paraphrased  & \textbf{0.927} & \textbf{0.891} \\
\midrule
\multirow{2}{*}{1 - TER} & Original & 0.274 & 0.760 \\
                        %\cline{2-4}
                         & Paraphrased  & 0.283 & 0.781 \\
\bottomrule
\end{tabular}

\caption{Pearson's correlation between automatic evaluation metrics and the silver standard of human judgments. The Paraphrased references are created using the best performing method from the previous section -- \emph{one-word only} with \emph{full Depfix}. To ensure consistency across metrics where higher values indicate better performance, we report 1−TER instead of the raw error rate.}
\label{choosing}
\end{table}

As a first step, we compare the effect of using paraphrased references on three widely used metrics: \bleu, \mn, and \ter. The results presented in \Cref{choosing} show that paraphrased references consistently improve correlation with human judgments for \bleu, \ter, and \mn.  
This improvement is somewhat unexpected, given that \mn internally integrates the same paraphrase table (\cmpt), whose subset (\fm) is used to generate our paraphrased references.

We therefore formulate the following hypothesis: while paraphrases are useful for bringing reference sentences closer to valid system outputs, using noisy paraphrase tables directly inside the evaluation metric may reduce its ability to discriminate between systems.  Importantly, in our approach, paraphrased references are subsequently corrected using Depfix, which can resolve grammatical inconsistencies introduced by substitution. In contrast, when paraphrase tables are used internally by the metric (as in \mn), no such correction is applied, and noisy or ill-formed paraphrases may be rewarded directly. This asymmetry suggests that paraphrases are more effective when used to generate alternative, post-edited reference formulations prior to evaluation rather than being incorporated directly into the metric.

In this section,\footnote{This section is based on \cite{parmesan}.} we explore an alternative approach in which paraphrases are used to generate additional reference variants rather than being handled internally by the metric. This leads to the formulation of the \parmesan metric described later in the section.


\subsection{Experimental Setup: Metric Configurations}

To test whether internal paraphrase matching is suboptimal, we evaluate \mn under several configurations that progressively reduce or entirely remove the internal paraphrase tables. These configurations are summarized in \Cref{parmesan:settings}.

\begin{table}[tb]
\centering
\resizebox{\linewidth}{!}{
\begin{tabular}{lrl}
\toprule
\mn setting & size & paraphrase table description \\
\midrule
\textbf{Original} & 684k & \cmpt \\
\textbf{One word} & 181k & only one-word pairs from \cmpt \\
\textbf{Same POS} & 122k & \textbf{One word}  with the same part-of-speech\\
\textbf{Different Lemma} & 71k & \textbf{Same POS} in which the pairs have different lemmas\\
\textbf{Same Lemma} & 51k & \textbf{Same POS} in which the pairs have the same lemma\\
\textbf{No paraphrases} & 0 & no paraphrase tables, i.e., exact match only\\
%\textbf{WordNet} & 202k & paraphrase tables generated from Czech WordNet\\
\bottomrule
\end{tabular}}
\caption{Different paraphrase tables for the \mn metric and their size (number of paraphrase pairs). The \textbf{Original} configuration corresponds to the original \mn metric, i.e., using the full paraphrase table without filtering.  Conversely, \textbf{No paraphrases} correspond to an empty set, so the \mn metric has to rely solely on exact match. 
% \ml{v pravém sloupci (popis tabulek) bych nepoužívala bold, mate to} 
%PB: vsude ty settings oznacuji boldem, tak musim i tyto. Prijde mi, ze vic mate, kdyz to tam nebude.
}
\label{parmesan:settings}
\end{table}

We test these \mn settings against three types of reference sentences:
\begin{description}
\item[Original:] original references provided by WMT for a particular year,
\item[Paraphrased:] references paraphrased by our \emph{One-word only} greedy substitution method using \fm and \wordnet, without any postprocessing.
\item[Depfixed:] \textbf{Paraphrased} references modified using \emph{full Depfix} for grammatical correction.
\end{description}

\subsection{Results on WMT12 and WMT13 Data}
\label{results-12-13}
We first report results on WMT12 and WMT13 in \Cref{tab:parmesan-results}. The results are consistent across both datasets. There is a notable positive impact of~the paraphrasing of~reference sentences and the application of Depfix. The strongest results are consistently obtained when paraphrases are used exclusively to generate paraphrased references, while the evaluation itself relies on exact matching.

\begin{table}[htb!]
\centering
\begin{tabular}{l|ccc}
%\multicolumn{7}{c}{One-word Only}\\

\multicolumn{4}{c}{\textbf{WMT12, silver standard}}  \\
\toprule
 & \multicolumn{3}{c}{\textbf{References}} \\
\mn setting & Original  & Paraphrased  & Depfixed \\
\midrule
\textbf{Original}     & \underline{0.833} & 0.905 & 0.927 \\
\textbf{One-word}      & 0.836 & 0.908 & 0.930  \\
\textbf{Same POS}      & 0.840 & 0.911 & 0.931  \\
\textbf{Same Lemma}     & 0.838 & 0.911 & 0.932 \\
\textbf{Different Lemma}    & 0.863 & 0.931 & 0.950  \\
\textbf{No paraphrases}     & 0.861 & 0.931 & \boxed{\textbf{0.951}}  \\
%\textbf{WordNet }        & 0.863 & 0.931 & \textbf{0.951} \\
\bottomrule
\end{tabular} 

\vspace{10pt}

\begin{tabular}{l|ccc}
%\multicolumn{7}{c}{One-word Only}\\
\multicolumn{4}{c}{\textbf{WMT13, silver standard} } \\
\toprule
 & \multicolumn{3}{c}{\textbf{References}} \\
\mn setting  & Original  & Paraphrased  & Depfixed\\
\midrule
\textbf{Original}        & \underline{0.817} & 0.865 & 0.891\\
\textbf{One-word }       & 0.820 & 0.867 & 0.892  \\
\textbf{Same POS}        & 0.823 & 0.869 & 0.893 \\
\textbf{Same Lemma }     & 0.821 & 0.868 & 0.892 \\
\textbf{Different Lemma }& 0.850 & 0.895  & \textbf{0.915}\\
\textbf{No paraphrases}  & 0.848 & 0.895 & \boxed{\textbf{0.915}} \\
%\textbf{WordNet }        & 0.850 & 0.894 & \textbf{0.915}\\
\bottomrule
\end{tabular}

\vspace{10pt}

\begin{tabular}{l|ccc}
 \multicolumn{4}{c}{\textbf{WMT13, gold standard}} \\
\toprule
 & \multicolumn{3}{c}{\textbf{References}} \\
\mn setting & Original  & Paraphrased & Depfixed\\
\midrule
\textbf{Original}           & \underline{0.856} & 0.894 & 0.918 \\
\textbf{One-word}        & 0.859 & 0.896 & 0.918 \\
\textbf{Same POS}        & 0.862 & 0.898 & 0.919 \\
\textbf{Same Lemma}      & 0.860 & 0.897 & 0.919 \\
\textbf{Different Lemma} & 0.885 & 0.918 & \textbf{0.933}\\
\textbf{No paraphrases}  & 0.883 & 0.917 & \boxed{\textbf{0.933}} \\
%\textbf{WordNet}         & 0.884 & 0.917 &  \textbf{0.933}\\
\bottomrule
\end{tabular}
\caption{Pearson's correlation between \mn and human judgments. For WMT12, the gold standard results are identical to the silver standard.} 
%korelace metrik - 0.981
\label{tab:parmesan-results}
\end{table}

The original \mn metric (i.e. the \textbf{Original} \mn setting with original references, underscored) shows the weakest performance for both WMT12 and WMT13. 

Filtering the paraphrase tables used within \mn consistently improves its performance, regardless of the reference sentence type. On the original reference sentences, \textbf{Different Lemma} setting yields the best results, as it effectively corresponds to the \fm tables.

We hypothesize that the main reason for the poorer performance of \mn when using \cmpt is the high level of noise in the paraphrase table. A substantial portion of this noise consists of pairs  corresponding to different inflectional forms of the same lemma. Such pairs can lead the metric to reward ungrammatical outputs, which is reflected in the low correlation observed for the \textbf{Same Lemma} setting.

An additional issue is that the metric may reward untranslated segments, as \cmpt contains English words and their Czech translations as paraphrase pairs (see \Cref{meteori}). As a result, differences between systems become less distinct, and the metric performs worse than when paraphrase matching is not used.

We also perform a control experiment in which reference sentences are paraphrased using the \emph{one-word only} method with word pairs from \cmpt (instead of \fm and \wordnet). Despite being much larger, \cmpt results in roughly half as many substitutions because it is not lemmatized (see \Cref{substitution_frequency}). The resulting correlations with the silver standard (0.869 and 0.852 for WMT12 and WMT13) are lower than those obtained using paraphrases from \fm and \wordnet (See the Paraphrased column in \Cref{tab:parmesan-results}). This confirms that the quality of the paraphrase tables plays an important role when generating alternative reference formulations.

\begin{table}[tb]
\centering
\begin{tabular}{cccc}
\toprule
paraphrases & \textbf{WMT12} & \textbf{WMT13} \\
\midrule
\wordnet & 0.26 & 0.22  \\
\fm &  1.53 & 1.29  \\
\wordnet $+$ \fm & 1.59 & 1.34  \\
\midrule
\cmpt &  0.75 & 0.61 \\
\bottomrule
\end{tabular}
\caption{Average number of one-word paraphrases per sentence.} 
\label{substitution_frequency}
\end{table}


\subsection{Defining the \parmesan Metric}
Based on our observations, we propose the \parmesan metric, our submission to the WMT14 Metrics Shared Task to evaluate English-to-Czech translation.

\parmesan (PARaphrasing for MEteor SANs paraphrases) operates in a two-stage process:
\begin{enumerate}
\item Reference sentences are paraphrased via the \emph{one-word only} greedy substitution method using \fm, followed by \emph{full Depfix}.
\item The resulting paraphrased references are evaluated using \mn without paraphrase matching, relying solely on exact matching.
\end{enumerate}

Unlike internal paraphrase matching, which applies uniformly, targeted paraphrasing operates asymmetrically and is guided by specific hypotheses, allowing it to introduce beneficial variation while avoiding excessive noise. In \Cref{tab:parmesan-results}, \parmesan scores are highlighted by enclosing them within a box, while the original \mn is underscored. The best scores are shown in bold. 

We use methods for comparing correlation coefficients \cite{meng1992comparing} to determine the statistical significance of~differences between metrics. The tests indicate that \parmesan outperforms \mn with 99\% certainty. 
%\ml{tenhle odstavec asi patří do předchozí sekce, příml k \Cref{tab:parmesan-results}??}
% PB: tam jeste neni definovany parmesan

%\scalebox{0.97}{
%\begin{tabular}{|l|ccc|ccc|}
%\hline
% & \multicolumn{3}{c|}{\textbf{WMT12} reference} &  \multicolumn{3}{c|}{\textbf{WMT13} reference} \\
%\hline
%\mn setting & Original  & Paraphr.  & Depfix.  & Original  & Paraphr.  & Depfix.\\
%\hline
%\textbf{Original}           & \underline{0.833} & 0.905 & 0.927 & \underline{0.817} & 0.865 & 0.891\\
%\textbf{One-word }       & 0.836 & 0.908 & 0.930 & 0.820 & 0.867 & 0.892  \\
%\textbf{Same POS}        & 0.840 & 0.911 & 0.931 & 0.823 & 0.869 & 0.893 \\
%\textbf{Same Lemma }     & 0.838 & 0.911 & 0.932 & 0.821 & 0.868 & 0.892 \\
%\textbf{Different Lemma }    & 0.863 & 0.931 & 0.950 & 0.850 & 0.895  & \textbf{0.915}\\
%\textbf{No paraphrases}     & 0.861 & 0.931 & \boxed{\textbf{0.951}} & 0.848 & 0.895 & \boxed{\textbf{0.915}} \\
%\textbf{WordNet }        & 0.863 & 0.931 & \textbf{0.951} & 0.850 & 0.894 & \textbf{0.915}\\
%\hline
%\end{tabular}} 

\subsection{Results on WMT14 data}
\label{vysledky}

\begin{table}[tb!]
\centering

\begin{tabular}{lccc}
%\multicolumn{7}{c}{One-word Only}\\
\multicolumn{4}{c}{\textbf{WMT14, silver standard}} \\
\toprule
\mn setting & Original  & Paraphrased & Depfixed \\
\midrule
\textbf{Original}        & \underline{0.963} & 0.957 & 0.968 \\
\textbf{One-word}        & 0.967 & 0.958 & 0.965 \\
\textbf{Same POS}        & 0.965 & 0.959 & 0.969 \\
\textbf{Same Lemma}      & 0.968 & 0.959 & 0.969  \\
\textbf{Different Lemma} & 0.970 & 0.965 & 0.968 \\
\textbf{No paraphrases}  & \textbf{0.973} & 0.965 & \boxed{0.968} \\
%\textbf{WordNet}         & \textbf{0.973} & 0.965 & 0.968 \\
\bottomrule
\end{tabular}

\vspace{10pt}

\begin{tabular}{lccc}
%\multicolumn{7}{c}{One-word Only}\\
\multicolumn{4}{c}{\textbf{WMT14, gold standard}} \\
\toprule
\mn setting & Original  & Paraphrased & Depfixed \\
\midrule
\textbf{Original}        & \underline{0.967} & 0.958 & 0.966\\
\textbf{One-word}        & 0.968 & 0.959 & 0.966\\
\textbf{Same POS}        & 0.969 & 0.959 & 0.966\\
\textbf{Same Lemma}      & 0.972 & 0.960 & 0.967 \\
\textbf{Different Lemma} & 0.972 & 0.963 & 0.962\\
\textbf{No paraphrases}  & \textbf{0.974} &  0.963 & \boxed{0.962}\\
%\textbf{WordNet}         & \textbf{0.973} & 0.965 & 0.968 & \textbf{0.974} & 0.963  & 0.962 \\
\bottomrule
\end{tabular}

\caption{Pearson's correlation of \mn and the silver, respectively, the gold 
standard -- \textit{TrueSkill} \citep{NIPS2006_f44ee263}. Note that as opposed to official WMT14 
results \citep{wmt14}, the version 1.4 of \mn is still used in this table.}
\label{results14}
\end{table}

The WMT14 results, presented in \Cref{results14}, show distinct patterns compared to the preceding two years. \parmesan still outperforms \mn on the silver standard, but only by a marginal difference (0.005). However, when using the gold standard, \mn performs better by precisely the same margin. In particular, Pearson's correlation between the two standards remains exceptionally high at 0.997.

In contrast to previous years, paraphrased reference sentences are detrimental to the performance of \mn on the WMT14 dataset. Applying Depfix to paraphrased references is beneficial but does not outperform the original references. However, the observed differences are minimal (up to 0.012).

Similarly as in WMT12 and WMT13, reducing the paraphrase tables in \mn remains beneficial for both original and paraphrased reference sentences. The best results are achieved with the \textbf{No paraphrases} configuration on the original references, i.e., no paraphrases are used. 

Notably, the highest correlations (up to 0.973) are obtained without using any paraphrases at all, indicating that in this setting any form of paraphrasing does not provide benefits and may be even slightly detrimental. This observation is consistent with the official results of the WMT14 metrics task \citep{machacek-bojar-2014-results}, where simple string-matching metrics outperform more complex ones for English-to-Czech evaluation. In particular, \mn is outperformed by TER, NIST \citep{nist}, and CDER \citep{leusch-etal-2006-cder}.

\subsection{Effect of Data Composition and Ceiling Effects}
We examine factors that could explain the discrepancy in performance of the \parmesan metric observed in the WMT14 results compared to WMT12 and WMT13. A plausible explanation concerns the composition of the source-language data. In WMT12 and WMT13, the source sentences are drawn evenly from all competing languages;
meaning only a~fifth or sixth of the data, respectively, is originally written in English and the remaining $ \frac{4}{5} $, respectively, $\frac{5}{6} $ are translations from the other languages. 

In contrast, the WMT14 English–Czech dataset was constructed exclusively from sentences originally written in either English or Czech. These sentences therefore may exhibit different lexical variability than the mixed translated data used in previous years.

Translating sentences from another source language into English and subsequently into Czech might be a contributing factor to the beneficial effect of prior reference sentences paraphrasing. To test whether this difference explains the behavior observed on WMT14, we repeat the experiments on a restricted subset of WMT12 and WMT13 containing only sentences originally written in English or Czech. The results are shown in Table~\ref{limited}.

However, even under this restriction, \parmesan continues to outperform the original \mn metric, and the overall trends remain consistent with the results obtained on the full datasets. This suggests that the presence of translation chains in the source data does not explain the different behavior observed for WMT14.

\begin{table}[tb]
\centering

\begin{tabular}{lccc}
\multicolumn{4}{c}{\textbf{WMT12, silver standard, Cz and En original sentences}} \\
\toprule
Reference & Original  & Paraphrased  & Depfixed \\
\midrule
\textbf{Original}           & \underline{0.781} & 0.872 & 0.897 \\
\textbf{One-word}        & 0.779 & 0.872 & 0.897  \\
\textbf{Same POS}        & 0.782 & 0.874 & 0.897 \\
\textbf{Same Lemma}      & 0.772 & 0.874 & 0.897  \\
\textbf{Different Lemma} & 0.807 & 0.898 & \textbf{0.923} \\
\textbf{No paraphrases} & 0.798 & 0.899 & \boxed{\textbf{0.923}}  \\
%\textbf{WordNet}         & 0.801 & 0.899 & \textbf{0.923} & 0.844 & 0.877 & \textbf{0.902}\\
\bottomrule
\end{tabular}

\vspace{10pt}

\begin{tabular}{lccc}
\multicolumn{4}{c}{\textbf{WMT13, silver standard, Cz and En original sentences}} \\
\toprule
Reference & Original  & Paraphrased  & Depfixed \\
\midrule
\textbf{Original}        & \underline{0.805} & 0.843 & 0.874\\
\textbf{One-word}        & 0.810 & 0.846 & 0.877  \\
\textbf{Same POS}        & 0.813 & 0.849 & 0.878 \\
\textbf{Same Lemma}      & 0.813 & 0.848 & 0.877 \\
\textbf{Different Lemma} & 0.842 & 0.879  & 0.877\\
\textbf{No paraphrases}  & 0.840 & 0.877 & \boxed{\textbf{0.902}} \\
%\textbf{WordNet}         & 0.801 & 0.899 & \textbf{0.923} & 0.844 & 0.877 & \textbf{0.902}\\
\bottomrule
\end{tabular} 
 
\caption{Pearson's correlation of \mn and the silver standard on sentences originally Czech or English only. Human judgments are computed only on those sentences as well.}
\label{limited}
\end{table} 

The more likely explanation is a \emph{ceiling effect}. Correlation on WMT14 is already exceptionally high ($\approx 0.97$). In such a high-accuracy environment, the baseline metric leaves almost no room for improvement, and the introduction of any transformation (even high-quality paraphrasing) may introduce more statistical noise than signal.

\subsection{Conclusion}
In summary, we highlight the adverse impact of noise in~\cmpt on the performance of~the \mn metric. We demonstrate that substantial reduction of the paraphrase tables can be beneficial for targeted paraphrasing of reference sentences. The \mn scores computed without paraphrase support, when applied to these new reference sentences, exhibit a significantly improved correlation with human judgments compared to the original \mn scores, particularly on the WMT12 and WMT13 data. 

The results on the WMT14 data do not support the findings from the two previous years and the observed improvement  was marginal. Variations in data composition do not provide a sufficient explanation for the differing behavior observed on the WMT14 dataset. However, one important observation is that the correlation between automatic metrics and human judgments on WMT14 is already extremely high $(\approx 0.97)$. In such a setting, even small changes to the evaluation procedure may have little effect or may slightly reduce correlation due to statistical variance.

This suggests that the benefits of paraphrased reference sentences are most apparent in evaluation scenarios with substantial surface variation between hypotheses and references, where the baseline metric struggles to capture valid translation variation.



\begin{mdframed}[
    linecolor=gray,
    linewidth=2pt,
    roundcorner=5,
    backgroundcolor=gray!5
]
\noindent 
We answer the following research questions:

\paragraph{Does the improvement in correlation achieved by targeted paraphrasing of reference sentences for \bleu also generalize to other commonly used evaluation metrics, in particular \mn, which already incorporates paraphrase matching?}
Yes, but not universally. On WMT12 and WMT13, targeted paraphrasing of reference sentences significantly improves the correlation of \mn with human judgments. However, the improvement in correlation is not confirmed by the WMT14 results.

\paragraph{Does the inherent noise in \cmpt actually degrade the performance of the \mn metric itself?}
Yes, our results provide strong evidence that noise in \cmpt harms the performance of the \mn metric. Across multiple experiments, we observe that entirely removing paraphrase matching within \mn consistently leads to the highest correlations with human judgments. The original \mn configuration, which relies on the full paraphrase tables, has the weakest performance on WMT12 and WMT13. 

\end{mdframed}

